\section{Conclusiones}

El desarrollo del Sistema de Exámenes de Colocación para el Centro Boliviano Americano representa un avance significativo en la digitalización de los procesos académicos de la institución. A continuación se presentan las principales conclusiones:

\begin{itemize}
    \item La arquitectura Frontend + BaaS (Supabase) demostró ser una solución eficiente, eliminando la necesidad de un servidor REST intermedio y reduciendo la complejidad operativa del sistema.
    \item El diseño de la base de datos en Tercera Forma Normal (3FN) garantiza la integridad y consistencia de los datos, mientras que las funciones SQL y triggers automatizan las reglas de negocio críticas directamente en la base de datos.
    \item La implementación de Row Level Security (RLS) proporciona un control de acceso granular a nivel de fila, asegurando que cada usuario solo pueda acceder a los datos que le corresponden.
    \item El sistema de autenticación dual (CI o email) facilita el acceso a los estudiantes, eliminando la necesidad de recordar un correo electrónico institucional.
    \item La protección de resultados históricos mediante triggers SQL garantiza la integridad de las evaluaciones, cumpliendo con uno de los requisitos más críticos del sistema.
\end{itemize}

\section{Recomendaciones}

\begin{itemize}
    \item \textbf{Despliegue en producción}: Se recomienda configurar correctamente las variables de entorno en el entorno de producción, incluyendo las claves de Supabase y las URLs de la aplicación.
    \item \textbf{Carga de preguntas}: Para una operatividad completa, se recomienda que el administrador cargue un número suficiente de preguntas distribuidas equitativamente entre los niveles antes de poner el sistema en producción.
    \item \textbf{Migración de datos históricos}: Si el CBA cuenta con datos de evaluaciones anteriores, se recomienda diseñar un proceso de migración para incorporarlos al nuevo sistema.
    \item \textbf{Backup periódico}: Configurar backups automáticos de la base de datos en Supabase para garantizar la disponibilidad de la información.
    \item \textbf{Monitoreo}: Implementar monitoreo de rendimiento y alertas para detectar posibles cuellos de botella cuando el sistema alcance su capacidad máxima de 30 estudiantes simultáneos.
    \item \textbf{Ampliación futura}: El diseño modular del sistema permite incorporar nuevas funcionalidades en el futuro, como la reproducción de audio para exámenes de listening, integración con LMS institucional, o generación de certificados automáticos.
\end{itemize}
